The weaknesses of textbook RSA arise because the raw function
\[
M \mapsto M^e \bmod N
\]
is deterministic and algebraically structured. To convert RSA into a practical encryption scheme, one must first transform the message using a randomized encoding process. This transformation is called \textbf{padding}.

\subsection{Why Padding Is Needed}

Padding serves two essential purposes. First, it introduces \textbf{randomness}, ensuring that the same message does not always encrypt to the same ciphertext. Second, it adds \textbf{structure}, so that malformed or manipulated inputs can be recognized and rejected after decryption.

Conceptually, padding changes the encryption process from
\[
C = M^e \bmod N
\]
to
\[
C = \operatorname{RSA}\bigl(\operatorname{Pad}(M;r)\bigr) = \operatorname{Pad}(M;r)^e \bmod N,
\]
where $r$ is fresh randomness drawn independently for every encryption. Because the RSA permutation is now applied to a randomized encoding of the message rather than to the message itself, the determinism and the multiplicative structure exploited against textbook RSA no longer translate into useful attacks: equality tests, dictionary attacks, small-message root extraction, and chosen-ciphertext manipulation all become much harder.

A useful way to state the design goal is that padding should turn RSA into a scheme that is \textbf{semantically secure}, and ideally secure against \textbf{adaptive chosen-ciphertext attacks} (IND-CCA2), so that an adversary learns nothing from a ciphertext even when allowed to query a decryption oracle on other ciphertexts.

\subsection{PKCS\#1 v1.5}

One historically important padding scheme is \textbf{PKCS\#1 v1.5}. Here $k$ denotes the \emph{byte length of the modulus} $N$, i.e.\ the number of bytes needed to write $N$,
\[
k = \left\lceil \frac{\log_2 N}{8} \right\rceil ,
\]
so that any $k$-byte string can be read as an integer in the range $[0, N)$ and fed directly into the RSA operation. Before RSA encryption, a message $M$ of $\mathrm{mLen}$ bytes is embedded into an encoded block $\mathrm{EM}$ of exactly $k$ bytes,
\[
\mathrm{EM} = \underbrace{\mathtt{00}}_{1}\,\|\,\underbrace{\mathtt{02}}_{1}\,\|\,\underbrace{\text{PS}}_{\ge 8}\,\|\,\underbrace{\mathtt{00}}_{1}\,\|\,\underbrace{M}_{\mathrm{mLen}} ,
\qquad |\mathrm{EM}| = k ,
\]
where $\|$ denotes byte concatenation and each underbrace gives the length of that field in bytes.
The components have specific roles:

\begin{itemize}
\item the leading \texttt{00} guarantees that the block, read as an integer, is smaller than $N$;
\item the block type byte \texttt{02} signals that the block is intended for \emph{encryption};
\item the padding string PS consists of \textbf{nonzero} random bytes and must be at least $8$ bytes long, which is what supplies the randomness;
\item the \texttt{00} separator marks the end of the padding and the start of the message.
\end{itemize}

Because PS must occupy at least $8$ bytes and the fixed bytes occupy $3$ more, the message length is bounded by $\mathrm{mLen} \le k - 11$. To decrypt, one computes $C^d \bmod N$, re-expresses it as a $k$-byte block, and checks that it has the expected \texttt{00 02} prefix, a nonzero padding region, and a \texttt{00} separator; the message is whatever follows.

This design was widely deployed and was a major improvement over textbook RSA. However, it was later shown to be vulnerable to adaptive chosen-ciphertext attacks whenever an implementation revealed whether a decrypted block had valid PKCS\#1 formatting (see Section~7). PKCS\#1 v1.5 therefore illustrates a recurring theme in cryptography: an encoding rule can look strong on paper yet still fail when combined with an implementation that leaks even one bit about the decrypted value.

\subsection{OAEP}

A more modern and principled design is \textbf{Optimal Asymmetric Encryption Padding (OAEP)}, introduced by Bellare and Rogaway. OAEP preprocesses the message using fresh randomness together with two hash-based mask generation functions, denoted $G$ and $H$, before the RSA exponentiation is applied. At a high level it:

\begin{itemize}
\item adds a fresh random seed to every encryption,
\item mixes the message and the seed through a two-round Feistel-like network built from $G$ and $H$, and
\item produces an \emph{all-or-nothing} encoded block: the message can be recovered only from the complete block, never from a part of it.
\end{itemize}

Concretely, let $\mathrm{hLen}$ be the output length of the hash and $k$ the byte length of $N$. The message is first formatted into a data block
\[
\mathrm{DB} = \mathrm{lHash} \,\|\, \mathrm{PS} \,\|\, \mathtt{01} \,\|\, M,
\]
where $\mathrm{lHash}$ is the hash of an optional label and $\mathrm{PS}$ is a run of zero bytes. A random seed of length $\mathrm{hLen}$ is then drawn, and the encoded message is computed as
\begin{align*}
\text{maskedDB} &= \mathrm{DB} \oplus G(\text{seed}), \\
\text{maskedSeed} &= \text{seed} \oplus H(\text{maskedDB}), \\
\mathrm{EM} &= \mathtt{00} \,\|\, \text{maskedSeed} \,\|\, \text{maskedDB},
\end{align*}
and $\mathrm{EM}$ is the integer fed into the RSA permutation. Decryption simply reverses the two rounds: $H$ recovers the seed from maskedSeed and maskedDB, then $G$ recovers DB, after which the formatting of DB is verified.

\begin{figure}[h]
\centering
\begin{tikzpicture}[>=Latex,node distance=0.6cm]
    \node[draw,rounded corners,fill=gray!10,minimum width=2.6cm,minimum height=0.9cm] (db)   at (0,4)   {DB ($\mathrm{lHash}\|\mathrm{PS}\|\mathtt{01}\|M$)};
    \node[draw,rounded corners,fill=gray!10,minimum width=2.0cm,minimum height=0.9cm] (seed) at (6.4,4) {seed (random)};

    \node[draw,circle,inner sep=1.5pt] (x1) at (0,2.4) {$\oplus$};
    \node[draw,rounded corners,fill=gray!10,minimum width=1.0cm,minimum height=0.8cm] (G) at (3.2,2.4) {$G$};

    \node[draw,rounded corners,fill=gray!10,minimum width=2.4cm,minimum height=0.9cm] (mdb) at (0,1.0) {maskedDB};

    \node[draw,rounded corners,fill=gray!10,minimum width=1.0cm,minimum height=0.8cm] (H) at (3.2,1.0) {$H$};
    \node[draw,circle,inner sep=1.5pt] (x2) at (6.4,1.0) {$\oplus$};

    \node[draw,rounded corners,fill=gray!10,minimum width=2.4cm,minimum height=0.9cm] (mseed) at (6.4,-0.4) {maskedSeed};

    \node[draw,rounded corners,fill=gray!12,minimum width=7.2cm,minimum height=0.9cm] (em) at (3.2,-1.9) {$\mathrm{EM}=\mathtt{00}\,\|\,\text{maskedSeed}\,\|\,\text{maskedDB}$};

    \draw[->,thick] (db.south) -- (x1.north);
    \draw[->,thick] (seed.south) .. controls (6.4,3.0) and (4.2,2.4) .. (G.east);
    \draw[->,thick] (G.west) -- (x1.east);
    \draw[->,thick] (x1.south) -- (mdb.north);
    \draw[->,thick] (mdb.east) -- (H.west);
    \draw[->,thick] (H.east) -- (x2.west);
    \draw[->,thick] (seed.east) .. controls (7.6,4) and (7.6,1.0) .. (x2.east);
    \draw[->,thick] (x2.south) -- (mseed.north);
    \draw[->,thick] (mdb.south) .. controls (0,-1.0) .. (em.north west);
    \draw[->,thick] (mseed.south) -- (em.north east);
\end{tikzpicture}
\caption{OAEP combines the data block DB and a random seed through two mask generation functions $G$ and $H$ before RSA is applied. Decryption must process the entire block: recovering $M$ requires all of maskedDB (to reconstruct the seed via $H$) and all of the seed (to reconstruct DB via $G$).}
\end{figure}

OAEP is important because it comes with stronger theoretical guarantees. The construction has the \emph{all-or-nothing} property: flipping a single bit of $\mathrm{EM}$ scrambles the recovered seed and data block unpredictably, so an attacker cannot transform a valid ciphertext into another valid one. In the random oracle model, RSA-OAEP is provably \textbf{IND-CCA2 secure} under the RSA assumption, which is precisely the resistance to adaptive chosen-ciphertext attacks that PKCS\#1 v1.5 lacked.

The broader lesson is that \textbf{secure RSA encryption is not just ``RSA plus a message''}. It is the RSA permutation composed with a carefully designed randomized encoding procedure. Without that preprocessing step, the trapdoor permutation remains mathematically elegant but cryptographically unsafe.
